\CUSTOMappendix{Дописывание получения движений токена по workflow}

\begin{lstlisting}[
  caption={Handler},
  label={lst:activity-log}
]
@Service
@RequiredArgsConstructor
public class WorkflowEventProcessor implements EventHandler {

  private static final Logger logger =
      LoggerFactory.getLogger(WorkflowEventProcessor.class);

  private final EventSender eventSender;
  private final ProcessObjectService processObjectService;
  private final ProcessObjectObserver processObjectObserver;
  private final ObjectProvider<StepUtilityService> stepUtility;

  @Override
  @Transactional
  public CompletableFuture<Void> onProcessInstance(ProcessEvent event) {
    return processObjectService.findByInstanceId(event.getRootInstanceId())
        .thenCompose(optional -> optional
            .map(obj -> {
              var holder = new AutoHolder<>(obj::getInput);
              var dirty = Holder.of(false);
              switch (obj.getStatus()) {
                case Status.FINISHED, Status.CANCELED:
                  break;
                default:
                  switch (event.getState()) {
                    case TERMINATED, EXTERNALLY_TERMINATED:
                      holder.getValue().setTransition(Status.CANCELED);
                      break;
                    case COMPLETED:
                      holder.getValue().setTransition(Status.FINISHED);
                      break;
                    default:
                      logger.info("id={} status={}", obj.getId(),
obj.getStatus());
                  }
                  if (obj.getEndDate() == null) {
                    holder.getValue().setEndDate(event.getEndTime());
                  }
                  dirty.setValue(true);
                  break;
              }
              if (holder.isPresent() && Boolean.TRUE.equals(dirty.getValue()))
{
                return processObjectObserver
                    .onFinished(new FinishedEvent(obj, holder.getValue(),
                        event.getState().name()))
                    .<Void>thenApply(r -> null);
              }
              return AsyncUtil.<Void>completedFuture(null);
            })
            .orElseGet(() -> AsyncUtil.runAsync(retry(event))));
  }

  @Override
  @Transactional
  public CompletableFuture<Void> onTaskInstance(TaskEvent event) {
    return processObjectService.findBaseByInstanceId(event.getRootInstanceId())
        .thenCompose(optional -> optional
            .<AsyncFuture<Void>>map(base -> {
              switch (event.getEventType()) {
                case DELETE, COMPLETE:
                  return stepUtility.getObject()
                      .closeTask(base, event.getTaskId());
                case CREATE:
                  return processObjectService
                      .loadFromRemote(base, Optional.of(event.getTaskKey()))
                      .thenCompose(obj -> handleCreate(event, obj));
                default:
                  logger.warn("skip {}", event);
                  return AsyncUtil.completedFuture(null);
              }
            }).orElseGet(() -> AsyncUtil.runAsync(retry(event))));
  }

  private AsyncFuture<Void> handleCreate(TaskEvent event, ProcessObject obj) {
    var template = obj.getType()
        .getStepTemplates().stream()
        .filter(t -> t.getKey().equals(event.getTaskKey()))
        .findAny();
    return stepUtility.getObject()
        .handleIncoming(event.getInstanceId(), obj,
            event.getTaskId(), event.getTaskKey(), template)
        .<Void>thenApply(r -> null);
  }

  @Override
  @Transactional
  public CompletableFuture<Void> onActivityInstance(ActivityEvent event) {
    if (!accept(event)) {
      return CompletableFuture.completedFuture(null);                          
    }
    return processObjectService.findByInstanceId(event.getInstanceId())        
        .thenCompose(optional -> optional
            .map((ProcessObject obj) -> {                                      
              if (Objects.equals(event.getActivityId(), obj.getStatus())) {
                return AsyncUtil.<Void>completedFuture(null);                  
              } 
              var state = prefix(event.getActivityType()) +                    
normalize(event.getActivityId());                                              
              if (state.equals(obj.getStatus())) {
                return AsyncUtil.<Void>completedFuture(null);                  
              } 
              var holder = new AutoHolder<>(obj::getInput);                    
              var lifecycle = obj.getType().getLifecycle();
              if (lifecycle.getState(obj.getStatus())                          
                  .getAllowedTransitions().contains(state)) {                  
                Optional.ofNullable(lifecycle.getTransition(state))            
                    .map(Transition::getDestination)                           
                    .ifPresentOrElse(
                        s -> holder.getValue().setTransition(state, s),        
                        () -> holder.getValue().setTransition(state));         
              }
              if (holder.isPresent()) {                                        
                return processObjectObserver.onPending(                        
                    new PendingEvent(obj, state, ""));
              }                                                                
              return AsyncUtil.<Void>completedFuture(null);
            })                                                                 
            .orElseGet(() -> AsyncUtil.runAsync(retry(event))));
  }

  private Runnable retry(BaseEvent event) {                                    
    return () -> {
      try {                                                                    
        Thread.sleep(100);
      } catch (InterruptedException e) {
        throw new RuntimeException(e);
      }                                                                        
      eventSender.send(event);
    };                                                                         
  }             

  private boolean accept(ActivityEvent event) {
    return switch (event.getActivityType()) {
      case "intermediateMessageCatch" -> true;
      default -> false;                                                        
    };
  }                                                                            
                
  private String prefix(String type) {
    return switch (type) {
      case "eventBasedGateway"         -> Prefixes.GATEWAY;
      case "intermediateMessageCatch"  -> Prefixes.CATCH;                      
      case "subProcess", "multiInstanceBody" -> Prefixes.SUBPROCESS;
      case "serviceTask"               -> Prefixes.SERVICE;                    
      default -> "";                                                           
    };                                                                         
  }                                                                            
                
  private String normalize(String id) {
    return id.contains("#") ? id.substring(0, id.indexOf("#")) : id;
  }                                                                            
}
\end{lstlisting}

\begin{lstlisting}[
  caption={BaseEvent},
  label={lst:activity-log}
]
@Getter
@SuperBuilder
@NoArgsConstructor
@ToString
public abstract class BaseEvent {
  private EventType eventType;
  private EventClass eventClass;
  private String rootProcessInstanceId;
  private String processInstanceId;
  private String processDefinitionKey;
  private Date endTime;
}
\end{lstlisting}

\begin{lstlisting}[
  caption={ActivityInstanceEvent},
  label={lst:activity-log}
]
@Getter
@SuperBuilder
@NoArgsConstructor
@ToString
public class ActivityInstanceEvent extends BaseEvent {
  private String activityId;
  private String activityType;
}
\end{lstlisting}

\begin{lstlisting}[
  caption={TaskInstanceEvent},
  label={lst:activity-log}
]
@Getter
@SuperBuilder
@NoArgsConstructor
@ToString
public class TaskInstanceEvent extends BaseEvent {
  private String taskId;
  private String taskDefinitionKey;
}
\end{lstlisting}

\begin{lstlisting}[
  caption={ProcessInstanceEvent},
  label={lst:activity-log}
]
@Getter
@SuperBuilder
@NoArgsConstructor
@ToString
public class ProcessInstanceEvent extends BaseEvent {
  private ProcessInstanceState state;
}
\end{lstlisting}

\begin{lstlisting}[
  caption={EventClass},
  label={lst:activity-log}
]
public enum EventClass {
  PROCESS_INSTANCE,
  ACTIVITY_INSTANCE,
  TASK_INSTANCE
}
\end{lstlisting}

\begin{lstlisting}[
  caption={EventType},
  label={lst:activity-log}
]
public enum EventType {
  START,
  END,
  CREATE,
  DELETE,
  COMPLETE;

  public static EventType find(String name) throws Exception {
    for (EventType eventType : values()) {
      if (name.equalsIgnoreCase(eventType.name().toLowerCase())) {
        return eventType;
      }
    }
    throw new Exception();
  }
}
\end{lstlisting}

\begin{lstlisting}[
  caption={ProcessInstanceState},
  label={lst:activity-log}
]

public enum ProcessInstanceState {
  ACTIVE,
  COMPLETED,
  TERMINATED,
  EXTERNALLY_TERMINATED,
  SUSPENDED,
  CANCELLED,
  NOT_DEFINED;

  public static ProcessInstanceState find(String name) throws Exception {
    for (ProcessInstanceState state : values()) {
      if (name.equalsIgnoreCase(state.name().toLowerCase())) {
        return state;
      }
    }
    throw new Exception();
  }
}
\end{lstlisting}

\begin{lstlisting}[
  caption={SystemEnvironmentUtil},
  label={lst:activity-log}
]
public class SystemEnvironmentUtil {

  public static String getEnvOrDefault(String environment, String def) {
    var env = System.getenv(environment);
    return env != null ? env : def;
  }
}
\end{lstlisting}

\begin{lstlisting}[
  caption={KafkaProducerConfig},
  label={lst:activity-log}
]

@Configuration
public class KafkaProducerConfig {

  public static final String EVENT_TOPIC_NAME = "camunda-events";

  private final String kafkaBrokers;
  private final String kafkaTopic;

  public KafkaProducerConfig() {
    kafkaBrokers = SystemEnvironmentUtil.getEnvOrDefault("EVENT_LISTENER_KAFKA_BROKERS",
        "kafka:9876");
    kafkaTopic = SystemEnvironmentUtil.getEnvOrDefault("EVENT_LISTENER_KAFKA_TOPIC",
        EVENT_TOPIC_NAME);
  }

  @Bean
  public JsonSerializer<BaseEvent> eventDtoSerializer(ObjectMapper objectMapper) {
    var serializer = new JsonSerializer<BaseEvent>(objectMapper);
    serializer.setAddTypeInfo(false);
    return serializer;
  }

  @Bean
  public ProducerFactory<String, BaseEvent> producerFactory(
      JsonSerializer<BaseEvent> eventDtoSerializer) {
    return new DefaultKafkaProducerFactory<>(
        Map.of(ProducerConfig.BOOTSTRAP_SERVERS_CONFIG, kafkaBrokers),
        new StringSerializer(),
        eventDtoSerializer

    );
  }

  @Bean
  public KafkaTemplate<String, BaseEvent> kafkaTemplate(
      ProducerFactory<String, BaseEvent> producerFactory) {
    return new KafkaTemplate<>(producerFactory);
  }

  @Bean
  public String kafkaTopicName() {
    return kafkaTopic;
  }
}
\end{lstlisting}

\begin{lstlisting}[
  caption={EventProducer},
  label={lst:activity-log}
]
@Service
@RequiredArgsConstructor
public class EventProducer {

  private final String kafkaTopicName;
  private final KafkaTemplate<String, BaseEvent> kafkaTemplate;

  public void send(BaseEvent eventDto) {
    kafkaTemplate.send(kafkaTopicName, eventDto);
  }
}
\end{lstlisting}

\begin{lstlisting}[
  caption={EventListenerPlugin},
  label={lst:activity-log}
]
@Setter
@SpringBootApplication
public class EventListenerPlugin implements ProcessEnginePlugin {

  @Override
  public void preInit(final ProcessEngineConfigurationImpl processEngineConfiguration) {
    var application = SpringApplication.run(EventListenerPlugin.class);
    var eventHandler = application.getBean(HistoryEventHandler.class);
    processEngineConfiguration.setHistoryEventHandler(eventHandler);
  }

  @Override
  public void postInit(final ProcessEngineConfigurationImpl processEngineConfiguration) {
  }

  @Override
  public void postProcessEngineBuild(final ProcessEngine processEngine) {}
}
\end{lstlisting}

\begin{lstlisting}[
  caption={EventHandlerKafka},
  label={lst:activity-log}
]
@Component
@RequiredArgsConstructor
public class EventHandlerKafka implements HistoryEventHandler {

  private final EventHandlerManager eventHandlerManager;
  private final EventProducer eventProducer;

  @Override
  public void handleEvent(final HistoryEvent event) {
    if (!eventHandlerManager.filterDefaultProcesses(event)) {
      eventHandlerManager
          .filterAndCreateHistoryEvent(event)
          .ifPresent(eventProducer::send);
    }
  }

  @Override
  public void handleEvents(final List<HistoryEvent> events) {
    events.forEach(this::handleEvent);
  }
}
\end{lstlisting}

\begin{lstlisting}[
  caption={EventHandlerManager},
  label={lst:activity-log}
]
@Component
@RequiredArgsConstructor
@Sl4j
public class EventHandlerManager {

  public boolean filterDefaultProcesses(HistoryEvent event) {
    if (event.getProcessDefinitionKey() == null) {
      return false;
    }

    return event.getProcessDefinitionKey().equals("invoice")
        || event.getProcessDefinitionKey().equals("ReviewInvoice");
  }

  public Optional<BaseEvent> filterAndCreateHistoryEvent(HistoryEvent event)
      throws Exception {
    if (event.getEventType() == null) {
      return Optional.empty();
    }

    if (event instanceof HistoricActivityInstanceEventEntity e) {
      return buildActivityInstanceEvent(e);
    } else if (event instanceof HistoricTaskInstanceEventEntity e) {
      return buildTaskInstanceEvent(e);
    } else if (event instanceof HistoricProcessInstanceEventEntity e) {
      return buildProcessInstanceEvent(e);
    } else {
      return Optional.empty();
    }
  }

  private Optional<BaseEvent> buildActivityInstanceEvent(
      HistoricActivityInstanceEventEntity event) throws Exception {
    if (event.getActivityId() == null
        || "userTask".equals(event.getActivityType())) {
      return Optional.empty();
    }

    return switch (event.getEventType()) {
      case "start" -> Optional.of(ActivityInstanceEvent.builder()
          .eventClass(EventClass.ACTIVITY_INSTANCE)
          .eventType(EventType.find(event.getEventType()))
          .rootProcessInstanceId(event.getRootProcessInstanceId())
          .processInstanceId(event.getProcessInstanceId())
          .processDefinitionKey(event.getProcessDefinitionKey())
          .activityId(event.getActivityId())
          .activityType(event.getActivityType())
          .build());
      default -> Optional.empty();
    };
  }

  private Optional<BaseEvent> buildTaskInstanceEvent(
      HistoricTaskInstanceEventEntity event) throws Exception {
    return switch (event.getEventType()) {
      case "delete", "complete", "create" -> Optional.of(TaskInstanceEvent.builder()
          .eventClass(EventClass.TASK_INSTANCE)
          .eventType(EventType.find(event.getEventType()))
          .rootProcessInstanceId(event.getRootProcessInstanceId())
          .processInstanceId(event.getProcessInstanceId())
          .processDefinitionKey(event.getProcessDefinitionKey())
          .taskDefinitionKey(event.getTaskDefinitionKey())
          .taskId(event.getTaskId())
          .build());
      default -> Optional.empty();
    };
  }

  private Optional<BaseEvent> buildProcessInstanceEvent(
      HistoricProcessInstanceEventEntity event) throws Exception {
    return switch (event.getEventType()) {
      case "update", "create", "start", "migrate" -> Optional.empty();
      default -> Optional.of(ProcessInstanceEvent.builder()
          .eventClass(EventClass.PROCESS_INSTANCE)
          .eventType(EventType.find(event.getEventType()))
          .state(ProcessInstanceState.find(event.getState()))
          .rootProcessInstanceId(event.getRootProcessInstanceId())
          .processInstanceId(event.getProcessInstanceId())
          .processDefinitionKey(event.getProcessDefinitionKey())
          .endTime(event.getEndTime())
          .build());
    };
  }
}
\end{lstlisting}

\begin{lstlisting}[
  caption={EventListenerHandler},
  label={lst:activity-log}
]
public interface EventListenerHandler {

  CompletableFuture<Void> consumeActivityInstanceEvent(ActivityInstanceEvent event);

  CompletableFuture<Void> consumeProcessInstanceEvent(ProcessInstanceEvent event);

  CompletableFuture<Void> consumeTaskInstanceEvent(TaskInstanceEvent event);
}
\end{lstlisting}

\begin{lstlisting}[
  caption={},
  label={lst:activity-log}
]
@Configuration
public class KafkaConsumerConfig {


  private final String kafkaBrokers;

  public KafkaConsumerConfig() {
    kafkaBrokers = SystemEnvironmentUtil.getEnvOrDefault("EVENT_LISTENER_KAFKA_BROKERS",
        "kafka:9876");
  }

  @Bean
  public ConsumerFactory<String, BaseEvent> beConsumerFactory(
      BaseEventDeserializer baseEventDeserializer) {
    return new DefaultKafkaConsumerFactory<>(
        Map.of(ConsumerConfig.BOOTSTRAP_SERVERS_CONFIG, kafkaBrokers,
            ConsumerConfig.GROUP_ID_CONFIG, "group",
            "enable.auto.commit", false),
        new StringDeserializer(),
        baseEventDeserializer
    );
  }

  @Bean
  public ConcurrentKafkaListenerContainerFactory<String, BaseEvent>
      baseEventListenerContainerFactory(ConsumerFactory<String, BaseEvent> beConsumerFactory) {
    var containerFactory = new ConcurrentKafkaListenerContainerFactory<String, BaseEvent>();
    containerFactory.setConcurrency(1);
    containerFactory.setConsumerFactory(beConsumerFactory);
    containerFactory.getContainerProperties()
        .setAckMode(ContainerProperties.AckMode.MANUAL_IMMEDIATE);
    return containerFactory;
  }
}
\end{lstlisting}

\begin{lstlisting}[
  caption={CamundaEventProcessMap},
  label={lst:activity-log}
]
@Component
public class CamundaEventProcessMap {
  private final Map<String, CamundaEventProcessDeque> map = new HashMap<>();

  public void put(BaseEvent event, Function<BaseEvent, CompletableFuture<Void>> processor) {
    synchronized (map) {
      map.computeIfAbsent(event.getRootProcessInstanceId(),
          id -> new CamundaEventProcessDeque(map, processor, () -> map.remove(id)))
          .add(event);
    }
  }
}
\end{lstlisting}

\begin{lstlisting}[
  caption={CamundaEventProcessDeque},
  label={lst:activity-log}
]
public class CamundaEventProcessDeque {

  private final Queue<BaseEvent> queue = new ArrayDeque<>();
  private Object lock;
  private Function<BaseEvent, CompletableFuture<Void>> processor;
  private Runnable remove;

  public CamundaEventProcessDeque(Object lock,
      Function<BaseEvent, CompletableFuture<Void>> processor, Runnable remove) {
    this.lock = lock;
    this.processor = processor;
    this.remove = remove;

    CompletableFuture.runAsync(this::processLoop);
  }
  public void add(BaseEvent event) {
    queue.add(event);
  }

  private void processLoop() {
    for (;;) {
      BaseEvent currentEvent;
      synchronized (lock) {
        currentEvent = queue.poll();
        if (Objects.isNull(currentEvent)) {
          remove.run();
          break;
        }
      }
      processor.apply(currentEvent).<Void>exceptionally(e -> {
        LOG.warn(currentEvent.toString(), e);
        return null;
      }).join();
    }
  }
}
\end{lstlisting}

\begin{lstlisting}[
  caption={bpm-platform.xml },
  label={lst:activity-log}
]
<?xml version="1.0" encoding="UTF-8"?>
<bpm-platform xmlns="http://www.camunda.org/schema/1.0/BpmPlatform" xmlns:xsi="http://www.w3.org/2001/XMLSchema-instance"
xsi:schemaLocation="http://www.camunda.org/schema/1.0/BpmPlatform http://www.camunda.org/schema/1.0/BpmPlatform ">

<process-engine name="default">
  <job-acquisition>default</job-acquisition>
  <configuration>org.camunda.bpm.engine.impl.cfg.StandaloneProcessEngineConfiguration</configuration>
  <datasource>java:jdbc/ProcessEngine</datasource>

  <properties>
    <property name="history">full</property>
  </properties>

  <plugins>

    <plugin>
      <class>ru.rabis.camunda.eventlistener.EventListenerPlugin</class>
    </plugin>
  </plugins>


</process-engine>

</bpm-platform>
\end{lstlisting}

\begin{lstlisting}[
  caption={EventListenerConsumer},
  label={lst:activity-log}
]
@Service
@RequiredArgsConstructor
public class EventListenerConsumer {

  private final EventListenerHandler eventListenerHandler;
  private final CamundaEventProcessMap map;

  @KafkaListener(topics = KafkaConsumerConfig.EVENT_TOPIC_NAME,
      containerFactory = "baseEventListenerContainerFactory")
  public void consumeEvent(BaseEvent event, Acknowledgment acknowledgment) {
    map.put(event, e -> {
      if (e instanceof ActivityInstanceEvent e2) {
        return eventListenerHandler.consumeActivityInstanceEvent(e2);
      } else if (e instanceof ProcessInstanceEvent e2) {
        return eventListenerHandler.consumeProcessInstanceEvent(e2);
      } else if (e instanceof TaskInstanceEvent e2) {
        return eventListenerHandler.consumeTaskInstanceEvent(e2);
      }
      return CompletableFuture.failedFuture(new ClassCastException(e.getClass().getSimpleName()));
    });
    acknowledgment.acknowledge();
  }
}
\end{lstlisting}